% Part: second-order-logic
% Chapter: syntax-and-semantics
% Section: introduction

\documentclass[../../../include/open-logic-section]{subfiles}

\begin{document}

\olfileid{sol}{syn}{int}

\olsection{ভূমিকা}

প্রথম-ক্রমের যুক্তিবিদ্যায় কোনো প্রদত্ত ভাষার অ-যৌক্তিক সংকেতগুলি,
অর্থাৎ তার !!{constant}s, !!{function}s ও !!{predicate}s-কে যৌক্তিক
সংকেতগুলির সঙ্গে মিলিয়ে প্রথম-ক্রমীয় !!{structure}s সম্বন্ধে কথা
প্রকাশ করা হয়। কাজটি করা হয় পরিতৃপ্তির ধারণার সাহায্যে। এই ধারণা
!!a{structure}~$\Struct{M}$, একটি চলরাশি-আরোপ~$s$ এবং
!!a{formula}~$!A$-কে পরস্পরের সঙ্গে যুক্ত করে: $\Sat{M}{!A}[s]$ সিদ্ধ
হয় যদি এবং কেবল যদি তার !!{constant}s, !!{function}s ও
!!{predicate}s-কে $\Struct{M}$-এর নির্দেশমতো এবং তার মুক্ত চলরাশিগুলিকে
$s$-এর নির্দেশমতো ব্যাখ্যা করলে $!A$ যা প্রকাশ করে তা সত্য হয়।
!!{identity}~$\eq$-এর ব্যাখ্যা যেমন $\Sat{M}{!A}[s]$-এর সংজ্ঞার মধ্যেই
নির্মিত, $\lforall$ ও~$\lexists$-এর ব্যাখ্যাও তেমনই। প্রথমটি সব সময়
গঠনের !!{domain}~$\Domain{M}$-এর উপর অভেদ সম্পর্ক হিসেবে ব্যাখ্যাত হয়,
আর পরিমাণসূচকগুলি সব সময় গোটা !!{domain}-এর উপর বিস্তৃত হয়। কিন্তু
নির্ণায়ক কথাটি হলো, পরিমাণায়ন শুধু !!{domain}-এর উপাদানের উপর করা
যায়; তাই পরিমাণসূচকের পরে কেবল বস্তু-!!{variable}s বসতে পারে।

দ্বিতীয়-ক্রমের যুক্তিবিদ্যায় ভাষা ও পরিতৃপ্তির সংজ্ঞা উভয়কেই প্রসারিত
করে মুক্ত ও বদ্ধ অপেক্ষক-চলরাশি ও বিধেয়-চলরাশি এবং তাদের উপর পরিমাণায়ন
অন্তর্ভুক্ত করা হয়। বস্তু-চলরাশির সঙ্গে !!{constant}s-এর যে সম্পর্ক,
এই চলরাশিগুলির সঙ্গে যথাক্রমে !!{function}s ও !!{predicate}s-এরও সেই
সম্পর্ক। দ্বিতীয়-ক্রমীয় পদ ও !!{formula}s গঠনে এরা একই ভূমিকা পালন
করে, এবং এদের উপর পরিমাণায়নও অনুরূপভাবে সামলানো হয়। \emph{মানক}
অর্থতত্ত্বে দ্বিতীয়-ক্রমীয় পরিমাণসূচকগুলি উপযুক্ত ধরনের সম্ভাব্য সব
বস্তুর উপর বিস্তৃত হয় (অপেক্ষক-চলরাশির ক্ষেত্রে $\Domain{M}$ থেকে
$\Domain{M}$-এ $n$-স্থানীয় অপেক্ষক, আর বিধেয়-চলরাশির ক্ষেত্রে
$n$-স্থানীয় সম্পর্ক)। যেমন,
$\lforall[\Obj{v_0}][(\Obj{P^1_0}(\Obj{v_0}) \lor \lnot
  \Obj{P^1_0}(\Obj{v_0}))]$ প্রথম-ক্রম ও দ্বিতীয়-ক্রম উভয়
যুক্তিবিদ্যাতেই একটি সূত্র; কিন্তু দ্বিতীয়টিতে আমরা আরও বিবেচনা করতে
পারি
$\lforall[\Obj{V^1_0}][\lforall[\Obj{v_0}][(\Obj{V^1_0}(\Obj{v_0}) \lor
    \lnot \Obj{V^1_0}(\Obj{v_0}))]]$ এবং
$\lexists[\Obj{V^1_0}][\lforall[\Obj{v_0}][(\Obj{V^1_0}(\Obj{v_0}) \lor
    \lnot \Obj{V^1_0}(\Obj{v_0}))]]$। এগুলিতে কোনো মুক্ত চলরাশি নেই,
তাই এগুলি দ্বিতীয়-ক্রমের যুক্তিবিদ্যার !!{sentence}s। এখানে
$\Obj{V^1_0}$ একটি দ্বিতীয়-ক্রমীয় $1$-স্থানীয় বিধেয়-চলরাশি।
$\Obj{V^1_0}$-এর গ্রহণযোগ্য ব্যাখ্যাগুলি ঠিক সেগুলিই, যেগুলি
$\Obj{P^1_0}$-এর মতো কোনো $1$-স্থানীয় !!{predicate}-কে আরোপ করা যায়;
অর্থাৎ, $\Domain{M}$-এর উপসেটগুলি। তাই এদের উপর পরিমাণায়নের অর্থ,
$\lforall[\Obj{v_0}][(\Obj{V^1_0}(\Obj{v_0}) \lor \lnot
  \Obj{V^1_0}(\Obj{v_0}))]$-টি $\Obj{V^1_0}$-এর মান হিসেবে
$\Domain{M}$-এর কোনো উপসেট আরোপের সব উপায়ের জন্য, অথবা অন্তত একটির
জন্য, সিদ্ধ হয়। প্রত্যেক সেটে কোনো প্রদত্ত বস্তু হয় থাকে, নয় থাকে
না; তাই যেকোনো !!{structure}-এ উভয় বাক্যই সত্য।
% BN-SRC-253 TARGET-CORRECTION-BEGIN
\par\smallskip\noindent\textnormal{\small\textbf{উৎস-সংশোধন (BN-SRC-253):} হিমায়িত ইংরেজি উৎসে শেষ প্রদর্শিত সূত্রের এক ঘটনায় বস্তু-চলরাশি $\Obj v_0$ এবং অন্য ঘটনায় কাঁচা $v_0$ লেখা ছিল, যদিও আগেই ঘোষিত সূত্রে উভয় স্থানেই $\Obj{v_0}$ আছে। বাংলা পাঠে দুই ঘটনাতেই ঘোষিত বস্তু-চলরাশির লিপি পুনরুদ্ধার করা হয়েছে।} % BN-SRC-253 TARGET-CORRECTION-END
\end{document}
